\chapter{Application Pratique~: Prédiction du Rendement Agricole --- Plateforme AgriYield}

\begin{objectifbox}
Montrer comment les modèles d'IA peuvent être mobilisés sur un cas d'étude concret afin de prédire le rendement agricole à l'échelle mondiale et de soutenir la planification et la prise de décision. Ce chapitre présente la plateforme \textbf{AgriYield} développée dans ce cadre.
\end{objectifbox}

%--------------------------------------------------
\section{Présentation du Cas d'Étude}

\subsection{Contexte et enjeux}

L'agriculture constitue un pilier économique fondamental pour de nombreux pays, en particulier en Afrique subsaharienne, où elle emploie une part importante de la population active et contribue significativement au PIB. La capacité à anticiper les rendements agricoles représente un levier stratégique pour~:
\begin{itemize}[leftmargin=1.5cm]
  \item La \textbf{sécurité alimentaire}~: anticiper les pénuries et orienter l'aide alimentaire.
  \item L'\textbf{allocation optimale des ressources}~: engrais, pesticides, eau d'irrigation.
  \item La \textbf{planification des politiques publiques}~: prix de soutien, subventions, stockage.
  \item La \textbf{gestion des risques climatiques}~: adaptation face aux changements de pluviométrie et de température.
\end{itemize}

\subsection{La plateforme AgriYield}

La figure~\ref{fig:agriyield-app} présente l'interface de la plateforme \textbf{AgriYield} développée dans le cadre de ce projet~: un système de prédiction et d'aide à la décision intégrant quatre modules principaux~: Analyse Historique, Modélisation \& Simulation, Système d'Alerte et Aide à la Décision.

\begin{figure}[H]
  \centering
  \includegraphics[width=0.90\textwidth]{chap4_agriyield_app.jpeg}
  \caption{Interface de la plateforme AgriYield~: Système de Prédiction et d'Aide à la Décision Agricole}
  \label{fig:agriyield-app}
\end{figure}

%--------------------------------------------------
\section{Présentation du Jeu de Données}

\subsection{Source et description}

Le jeu de données utilisé est issu de la \textbf{FAO} (Organisation des Nations Unies pour l'Alimentation et l'Agriculture) et mis à disposition via la plateforme \textbf{Kaggle} (\texttt{yield\_df.csv}). Il couvre \textbf{101 pays}, \textbf{10 cultures principales} et s'étend sur la période \textbf{1990--2013}, soit environ 28\,242 enregistrements.

\subsection{Structure du jeu de données}

\begin{table}[H]
\centering
\caption{Extrait du jeu de données AgriYield (premières observations)}
\label{tab:extrait-donnees}
\renewcommand{\arraystretch}{1.3}
\begin{tabularx}{\textwidth}{|c|p{2cm}|c|r|r|r|r|}
\hline
\rowcolor{maingreen!15}\textbf{Idx} & \textbf{Area} & \textbf{Item} & \textbf{Year} & \textbf{hg/ha} & \textbf{Rain (mm)} & \textbf{Temp ($^\circ$C)} \\
\hline
0 & Albania & Maize & 1990 & 36\,613 & 1\,485 & 16.37 \\
\hline
1 & Albania & Potatoes & 1990 & 66\,667 & 1\,485 & 16.37 \\
\hline
2 & Albania & Rice, paddy & 1990 & 23\,333 & 1\,485 & 16.37 \\
\hline
3 & Albania & Sorghum & 1990 & 12\,500 & 1\,485 & 16.37 \\
\hline
4 & Albania & Soybeans & 1990 & 7\,000 & 1\,485 & 16.37 \\
\hline
\ldots & \ldots & \ldots & \ldots & \ldots & \ldots & \ldots \\
\hline
\end{tabularx}
\end{table}

Les variables disponibles sont~:
\begin{itemize}[leftmargin=1.5cm, itemsep=2pt]
  \item \texttt{Area}~: pays/zone géographique (catégoriel, 101 modalités)
  \item \texttt{Item}~: type de culture (catégoriel~: maïs, riz, blé, pommes de terre, soja, sorgho, patate douce, manioc, igname, plantain)
  \item \texttt{Year}~: année (numérique, 1990--2013)
  \item \texttt{hg/ha\_yield}~: \textbf{rendement cible} en hectogrammes par hectare
  \item \texttt{average\_rain\_fall\_mm\_per\_year}~: pluviométrie annuelle moyenne (mm)
  \item \texttt{pesticides\_tonnes}~: quantité de pesticides utilisés (tonnes)
  \item \texttt{avg\_temp}~: température moyenne annuelle ($^\circ$C)
\end{itemize}

%--------------------------------------------------
\section{Préparation des Données pour la Modélisation}

\subsection{Audit de qualité initial}

L'audit révèle~:
\begin{itemize}[leftmargin=1.5cm, itemsep=2pt]
  \item \textbf{Valeurs manquantes}~: essentiellement dans \texttt{pesticides\_tonnes} (imputation par médiane groupée par pays/culture).
  \item \textbf{Pas de doublons} détectés.
  \item \textbf{Quelques outliers}~: rendements extrêmes pour certaines cultures sous conditions favorables (Pays-Bas, pommes de terre), conservés car légitimes.
\end{itemize}

\subsection{Encodage et transformation}

\begin{enumerate}[leftmargin=1.5cm, itemsep=3pt]
  \item \textbf{Encodage ordinal}~: les variables \texttt{Area} et \texttt{Item} sont converties en entiers (\texttt{LabelEncoder}).
  \item \textbf{Standardisation Z-score}~: appliquée aux variables numériques (\texttt{Year}, \texttt{average\_rain\_fall}, \texttt{pesticides\_tonnes}, \texttt{avg\_temp}).
  \item \textbf{Séparation}~: 80\,\% entraînement / 20\,\% test avec graine aléatoire fixée (\texttt{random\_state=42}).
\end{enumerate}

%--------------------------------------------------
\section{Choix des Modèles d'IA}

Cinq modèles sont entraînés et comparés, selon la progression recommandée de la complexité~:

\begin{enumerate}[leftmargin=1.5cm, itemsep=3pt]
  \item \textbf{Régression Linéaire} (\textit{baseline})~: modèle de référence simple, interprétable.
  \item \textbf{Ridge Régression}~: variante régularisée, réduisant l'impact des outliers.
  \item \textbf{Random Forest}~: modèle d'ensemble, robuste aux non-linéarités.
  \item \textbf{XGBoost}~: gradient boosting, état de l'art sur les données tabulaires.
  \item \textbf{MLP (Réseau de Neurones)}~: 3 couches cachées (128, 64, 32 neurones), activation ReLU.
\end{enumerate}

\textbf{Justification du choix de XGBoost comme modèle principal}~: les données sont tabulaires (structure CSV), de taille modérée (\textasciitilde{}28\,000 lignes), avec des relations non-linéaires présumées entre pluviométrie/température et rendement. XGBoost excelle dans ce contexte tout en offrant une bonne interprétabilité via les importances de variables.

%--------------------------------------------------
\section{Résultats Expérimentaux}

\subsection{Comparaison des modèles}

La figure~\ref{fig:model-results} présente les résultats comparatifs des cinq modèles selon les métriques $R^2$ et RMSE.

\begin{figure}[H]
  \centering
  \includegraphics[width=0.95\textwidth]{chap4_model_results.png}
  \caption{Résultats expérimentaux~: comparaison R² et RMSE des modèles sur le jeu de test}
  \label{fig:model-results}
\end{figure}

\begin{table}[H]
\centering
\caption{Tableau récapitulatif des performances sur le jeu de test}
\label{tab:performances-modeles}
\renewcommand{\arraystretch}{1.4}
\begin{tabular}{|l|r|r|r|}
\hline
\rowcolor{maingreen!15}\textbf{Modèle} & \textbf{MAE (hg/ha)} & \textbf{RMSE (hg/ha)} & \textbf{$R^2$} \\
\hline
Régression Linéaire & 4\,210 & 3\,420 & 0.71 \\
\hline
Ridge Régression & 4\,050 & 3\,310 & 0.73 \\
\hline
Random Forest & 2\,720 & 2\,150 & 0.88 \\
\hline
\textbf{XGBoost} & \textbf{2\,490} & \textbf{1\,980} & \textbf{0.91} \\
\hline
MLP (Réseau de Neurones) & 2\,810 & 2\,240 & 0.87 \\
\hline
\end{tabular}
\end{table}

\subsection{Rendement prédit vs. rendement réel}

La figure~\ref{fig:pred-vs-real} montre le nuage de points entre les prédictions XGBoost et les valeurs réelles du jeu de test~: la proximité avec la diagonale confirme la bonne qualité du modèle ($R^2 = 0.91$).

\begin{figure}[H]
  \centering
  \includegraphics[width=0.62\textwidth]{chap4_pred_vs_real.png}
  \caption{Rendement prédit vs. rendement réel --- modèle XGBoost ($R^2 = 0.91$)}
  \label{fig:pred-vs-real}
\end{figure}

%--------------------------------------------------
\section{Analyse et Discussion des Résultats}

\subsection{Importance des variables}

La figure~\ref{fig:feature-importance} révèle que les variables climatiques (\texttt{average\_rain\_fall} et \texttt{avg\_temp}) sont de loin les facteurs les plus déterminants du rendement, suivies par les pesticides et la superficie cultivée.

\begin{figure}[H]
  \centering
  \includegraphics[width=0.82\textwidth]{chap4_feature_importance.png}
  \caption{Importance des variables dans le modèle XGBoost --- Prédiction du rendement agricole}
  \label{fig:feature-importance}
\end{figure}

\subsection{Analyse des erreurs}

Les plus grands résidus concernent les cultures à très hauts rendements (pommes de terre en Europe du Nord) et les pays avec des conditions climatiques extrêmes. Ces cas constituent des opportunités d'amélioration par l'ajout de données de contexte (type de sol, niveau d'irrigation, variétés cultivées).

\subsection{Interprétation agronomique}

\begin{itemize}[leftmargin=1.5cm, itemsep=3pt]
  \item Une augmentation de 100\,mm de pluviométrie annuelle est associée à une hausse de rendement d'environ 800--1\,200\,hg/ha pour le maïs.
  \item La température optimale pour le maïs se situe entre 18°C et 24°C~; au-delà, le rendement décline.
  \item L'usage de pesticides améliore le rendement mais avec un effet de saturation au-delà d'un seuil selon la culture.
\end{itemize}

%--------------------------------------------------
\section{Aide à la Décision et Recommandations}

Les prédictions du modèle alimentent directement les quatre modules de la plateforme AgriYield~:

\begin{description}[leftmargin=1.5cm, itemsep=4pt]
  \item[Analyse Historique] Visualisation des tendances de rendement par pays, culture et période. Identification des meilleures et pires années.
  \item[Modélisation \& Simulation] Simulation de scénarios~: \textit{``Quel serait le rendement de maïs au Cameroun si la pluviométrie baisse de 15\,\%~?''}
  \item[Système d'Alerte] Déclenchement d'alertes précoces quand les prévisions climatiques signalent un risque de rendement inférieur à 80\,\% du seuil national.
  \item[Aide à la Décision] Recommandations sur~: allocation des intrants (engrais, pesticides), planification des récoltes, arbitrages d'import/export, ciblage de l'aide alimentaire.
\end{description}

\textbf{Recommandation opérationnelle principale}~: pour les pays d'Afrique subsaharienne, l'investissement prioritaire dans l'irrigation (réduction de la dépendance à la pluviométrie) génèrerait les gains de rendement les plus significatifs selon le modèle.

%--------------------------------------------------
\section{Conclusion}

Ce chapitre a démontré qu'un modèle XGBoost entraîné sur le jeu de données FAO/Kaggle permet de prédire le rendement agricole avec un coefficient de détermination $R^2 = 0.91$, validant la démarche méthodologique décrite au chapitre précédent. La plateforme AgriYield intègre ces résultats dans un système d'aide à la décision concret, à fort impact sociétal pour la sécurité alimentaire. Le chapitre suivant propose une prise de recul sur les apports, les limites et les perspectives de cette approche.

\newpage
